\section{It\^o Integral}

Let $M \in \mathcal{M}_0^2$. We begin with the simplest class of integrands.

If $H_t = Z \mathsf{1}_{(S,T]}(t),$ where $Z \in \mathcal{F}_S$ is bounded and $S \leq T$ are stopping times, define
$$
    \int_0^t H_s \diff M_s \equiv (H \cdot M)_t \equiv Z\bigl(M_{T \wedge t} - M_{S \wedge t}\bigr).
$$

We now show that $(H \cdot M)$ belongs to $\mathcal{M}_0^2$. As a first step, we prove that $(H \cdot M)$ is a martingale. By Lemma 1.6.2, it suffices to verify that for every stopping time $R$,
$$
\begin{cases}
    \mathbb{E}\abs{(H \cdot M)_R} < \infty,\\
    \mathbb{E}(H \cdot M)_R = 0.
\end{cases}
$$

\pf{
    Let $R$ be any stopping time. Since $Z$ is bounded, say $\abs{Z} \leq K$, we have
    $$
        \mathbb{E}\abs{(H \cdot M)_R}
        \leq K\mathbb{E}\abs{M_{T \wedge R}} + K\mathbb{E}\abs{M_{S \wedge R}}.
    $$
    Because $M \in \mathcal{M}_0^2$, the martingale $M$ is uniformly integrable, and
    $$
        M_{T \wedge R} = \mathbb{E}(M_\infty \mid \mathcal{F}_{T \wedge R}).
    $$
    Hence
    $$
        \mathbb{E}\abs{M_{T \wedge R}}
        \leq \mathbb{E}\left[\mathbb{E}(\abs{M_\infty} \mid \mathcal{F}_{T \wedge R})\right]
        = \mathbb{E}\abs{M_\infty},
    $$
    and similarly for $M_{S \wedge R}$. Therefore
    $$
        \mathbb{E}\abs{(H \cdot M)_R} < \infty.
    $$

    It remains to show that $\mathbb{E}(H \cdot M)_R = 0$. By linearity, it is enough to treat the case $Z = \mathsf{1}_A$ with $A \in \mathcal{F}_S$. Then
    $$
        \mathbb{E}(H \cdot M)_R
        = \mathbb{E}\left[\mathsf{1}_A(M_{T \wedge R} - M_{S \wedge R})\right].
    $$
    Define the stopping times
    $$
        S_A \equiv
        \begin{cases}
            S, & \text{on } A,\\
            \infty, & \text{on } A^c.
        \end{cases}
        \qquad
        T_A \equiv
        \begin{cases}
            T, & \text{on } A,\\
            \infty, & \text{on } A^c.
        \end{cases}
    $$
    Then
    $$
        \mathsf{1}_A M_{S \wedge R} = M_{S_A \wedge R},
        \qquad
        \mathsf{1}_A M_{T \wedge R} = M_{T_A \wedge R}.
    $$
    Hence the expectation above is the difference of expectations of stopped values of the martingale $M$, and therefore equals $0$ by optional sampling. Thus $(H \cdot M)$ is a martingale.
}

\lemnn{
    If $Z \in b\mathcal{F}_S$ and $M$ is a uniformly integrable martingale, then $H \cdot M$ is a uniformly integrable martingale.
}

More generally, consider a finite linear combination of such elementary processes:
$$
    H_t = \sum_{i=1}^n Z_i \mathsf{1}_{(S_i,T_i]}(t),
$$
where each $Z_i \in b\mathcal{F}_{S_i}$.

We define
$$
    (H \cdot M)_t
    \triangleq
    \sum_{i=1}^{n} Z_i \bigl(M_{T_i \wedge t} - M_{S_i \wedge t}\bigr).
$$

\cor{
    If $Z_i \in b \mathcal{F}_{S_i}$ and $M$ is U.I. martingale, then $(H\cdot M)_{t}$ is U.I. martingale.
}


We next introduce the predictable $\sigma$-algebra, which is the natural measurability class for It\^o integrands.

\defn{Previsible process}{
    The previsible $\sigma$-algebra $\mathcal{P}$ on $(0,\infty) \times \Omega$ is defined to be the smallest $\sigma$-algebra for which every adapted $L$-process is $\mathcal{P}$-measurable.

    An $L$-process is a process $(X_t)_{t>0}$ such that $X_t = X_{t-}$ and $X_{t+}$ exists.

    We say that a process $(X_t)_{t>0}$ is previsible if it is $\mathcal{P}$-measurable.
}

\lemnn{
    Suppose $S$ and $T$ are stopping times with $S \leq T \leq \infty$. Let $Z \in b\mathcal{F}_S$ and $H = Z\mathsf{1}_{(S,T]}$. Then $H$ is previsible.
}

\pf{
    Fix $t > 0$ and let $B$ be a Borel set. We may assume that $0 \notin B$. Then
    $$
        \{H_t \in B\}
        = \{Z \in B\} \cap \{S < t\} \cap \{t \leq T\}.
    $$
    Moreover,
    $$
        \{S < t\} = \bigcup_{n=1}^{\infty}\left\{S \leq t - \frac{1}{n}\right\}.
    $$
    Since $\{Z \in B\} \in \mathcal{F}_S$, we have
    $$
        \{Z \in B\} \cap \left\{S \leq t - \frac{1}{n}\right\} \in \mathcal{F}_{t - \frac{1}{n}} \subseteq \mathcal{F}_t.
    $$
    Also, $\{t \leq T\} \in \mathcal{F}_t$. Therefore $\{H_t \in B\} \in \mathcal{F}_t$, so $H$ is adapted. Since $H$ is clearly an $L$-process, it follows from the definition of $\mathcal{P}$ that $H$ is previsible.
}

\defnn{
    Let
    $$
        b\varepsilon
        = \left\{
            \sum_{i=1}^n Z_i \mathsf{1}_{(S_i,T_i]}
            :
            n \geq 1,\,
            Z_i \in b\mathcal{F}_{S_i},\,
            S_1 \leq T_1 \leq S_2 \leq \cdots
        \right\}.
    $$
}

\noindent\textbf{Claim.} $\sigma(b\varepsilon) = \mathcal{P}$.

\pf{
    ``$\subseteq$'' The inclusion $\sigma(b\varepsilon) \subseteq \mathcal{P}$ follows from the previous lemma.

    ``$\supseteq$'' For any bounded left-continuous process,
    $$
    H_t = \lim_{k \uparrow \infty} \lim_{n \uparrow \infty} \sum_{i=2}^{n_k} H_{\frac{i-1}{n}} \mathsf{1}_{\left(\frac{i-1}{n},\frac{i}{n}\right]}(t).
    $$

    Since $H$ is adapted, we get $H_{\frac{i-1}{n}} \in b\mathcal{F}_{\frac{i-1}{n}}$.

    Hence
    $$
    \sum_{i=2}^{n_k} H_{\frac{i-1}{n}} \mathsf{1}_{\left( \frac{i-1}{n}, \frac{i}{n} \right]}(t) \in b \varepsilon.
    $$  

    $\implies H_t \in \sigma(b \varepsilon).$

}
\pf{

    
}
\lemnn{
    Let $H \in b\varepsilon$ and $M \in \mathcal{M}_0^2$. Then $H \cdot M \in \mathcal{M}_0^2$, and if
    $$
        H = \sum_{i=1}^n Z_{i-1}\mathsf{1}_{(T_{i-1},T_i]},
    $$
    then
    $$
        \mathbb{E}(H \cdot M)_\infty^2
        =
        \sum_{i=1}^{n}\mathbb{E}\left[ Z_{i-1}^2(M_{T_i} - M_{T_{i-1}})^2 \right].
    $$
}


\pf{
    $$
        (H \cdot M)_t = \sum_{i=1}^n Z_{i-1}\left( M_{T_i \wedge t} - M_{T_{i-1} \wedge t} \right).
    $$
    By the previous corollary 1.7.1, $(H \cdot M)$ is a U.I. martingale.

    Hence $(H \cdot M)_t \to (H \cdot M)_\infty$ a.s. and in $L^1$, where
    $$
        (H \cdot M)_\infty = \sum_{i=1}^n Z_{i-1}\Delta_i,
        \qquad
        \Delta_i = M_{T_i} - M_{T_{i-1}}.
    $$

    Therefore
    $$
        \mathbb{E}\left[ (H \cdot M)_\infty^2 \right]
        = \mathbb{E} \sum_{i=1}^n Z_{i-1}^2 \Delta_i^2
        + 2 \mathbb{E}\left( \sum_{i < j} Z_{i-1}\Delta_i Z_{j-1}\Delta_j \right).
    $$

    For $i<j$, the random variable $Z_{i-1}\Delta_i Z_{j-1}$ is $\mathcal{F}_{T_{j-1}}$-measurable, so
    $$
        \mathbb{E}\left[ Z_{i-1}\Delta_i Z_{j-1}\Delta_j \right]
        =
        \mathbb{E}\left[
            Z_{i-1}\Delta_i Z_{j-1}
            \mathbb{E}\left( \Delta_j \middle| \mathcal{F}_{T_{j-1}} \right)
        \right]
        = 0.
    $$
    Thus all cross terms vanish.

    By Doob's $L^2$ inequality,
    $$
        \mathbb{E}\left[
            \sup_{t \geq 0}(H \cdot M)_t^2
        \right]
        \leq 4 \mathbb{E}\left[ (H \cdot M)_\infty^2 \right]
        < \infty.
    $$
    Therefore $H \cdot M \in \mathcal{M}_0^2$.
}

Recall from Meyer's theorem that for $M \in \mathcal{M}_0^2$ there exists a quadratic variation $[M]$ such that $M^2 - [M]$ is a uniformly integrable martingale. Applying this to the previous lemma 1.7.3, we get
$$
    \mathbb{E}(H \cdot M)_\infty^2
    =
    \sum_{i=1}^n \mathbb{E}\left[
        Z_{i-1}^2\bigl(M_{T_i} - M_{T_{i-1}}\bigr)^2
    \right].
$$
Moreover,
$$
    \begin{aligned}
        \mathbb{E}\left[
            \bigl(M_{T_i} - M_{T_{i-1}}\bigr)^2 \middle| \mathcal{F}_{T_{i-1}}
        \right]
        &=
        \mathbb{E}\left[
            M_{T_i}^2 - M_{T_{i-1}}^2 \middle| \mathcal{F}_{T_{i-1}}
        \right] \\
        &=
        \mathbb{E}\left[
            [M]_{T_i} - [M]_{T_{i-1}} \middle| \mathcal{F}_{T_{i-1}}
        \right],
    \end{aligned}
$$
because $M^2 - [M]$ is a martingale and $\mathbb{E}(M_{T_i}\mid \mathcal{F}_{T_{i-1}})=M_{T_{i-1}}$.
Therefore
$$
    \mathbb{E}(H \cdot M)_\infty^2
    =
    \sum_{i=1}^{n}\mathbb{E}\left[
        Z_{i-1}^2\bigl([M]_{T_i} - [M]_{T_{i-1}}\bigr)
    \right]
    =
    \mathbb{E}\left[
        \int_0^\infty H_s^2 \diff [M]_s
    \right].
$$

This shows that, for simple predictable integrands, the map
$$
    H \mapsto (H \cdot M)_\infty
$$
is an isometry from $(L^2(M)\cap b\varepsilon,\norm{\cdot}_M)$ into $L^2(\mathcal{F}_\infty)$ equipped with the usual $L^2$ norm.

For fixed $M \in \mathcal{M}_0^2$, define
$$
    \norm{H}_M
    =
    \left(
        \mathbb{E}\int_0^\infty H_s^2 \diff [M]_s
    \right)^{1/2}.
$$
Let
$$
    L^2(M)
    =
    \left\{
        \text{previsible } H : \norm{H}_M < \infty
    \right\}.
$$

Then the map
$$
\begin{aligned}
    I : L^2(M) \cap b\varepsilon &\to L^2(\mathcal{F}_\infty),\\
    H &\mapsto (H \cdot M)_\infty
\end{aligned}
$$
is an isometry, since
$$
    \mathbb{E}\left[
        \int_0^\infty H_s^2 \diff [M]_s
    \right]
    =
    \mathbb{E}(H \cdot M)_\infty^2.
$$
Hence $I$ extends uniquely from simple predictable processes to all of $L^2(M)$.

\defn{It\'o Integral}{
    Let $M \in \mathcal{M}_{0}^2$ and $H \in L^2(M)$. Then It\'o Integral $(H\cdot M)_{t}$ is the image of $H$ under extension of $I : L^2(M)\to L^2(\mathcal{F}_{\infty})$.
}

\noindent\textbf{Localization.}

\defnn{
    Let $T$ be a stopping time and let $H$ be any process. Define
    $$
        H(0,T](t)
        =
        \begin{cases}
            H_t, & 0 < t \leq T,\\
            0, & t > T.
        \end{cases}
    $$
}

\rmkb{
    If $H$ is previsible, then so is $H(0,T]$.
}

\defnn{
    Let $T$ be a stopping time and let $X$ be any process. Define
    $$
        X^T(t)
        =
        \begin{cases}
            X(t), & 0 \leq t \leq T,\\
            X(T), & t > T.
        \end{cases}
    $$
}
For an elementary integrand, the localization identity can be checked directly. Indeed,
$$
    (H \cdot M)^T_t = (H(0,T]\cdot M^T)_t,
$$
and
$$
    \left[ Z_i \left( M_{T_i \wedge t} - M_{T_{i-1} \wedge t} \right) \right]^T
    =
    \left( Z_i \mathsf{1}_{(T_{i-1}\wedge T,\; T_i \wedge T]} \cdot M^T \right)_t
    =
    Z_i\left( M_{T_i \wedge T \wedge t} - M_{T_{i-1}\wedge T \wedge t} \right).
$$
\thmnn{
    Fix $M \in \mathcal{M}_0^2$. For any $H \in L^2(M)$ and any stopping time $T$:
    \begin{itemize}
        \item[(i)] $(H \cdot M)^T = H(0,T] \cdot M = H \cdot M^T$.
        \item[(ii)] $(H \cdot M)_t^2 - \int_0^t H_s^2 \diff [M]_s$ is a uniformly integrable martingale. In particular,
        $$
            [H \cdot M]_t = \int_0^t H_s^2 \diff [M]_s.
        $$
        \item[(iii)] If $H, K \in b\mathcal{P}$, then
        $$
            (HK)\cdot M = H \cdot (K \cdot M) = K \cdot (H \cdot M).
        $$
        \item[(iv)] $\Delta(H \cdot M) = H \Delta M$.
        \item[(v)] If $M$ is continuous, then $H \cdot M$ is continuous.
    \end{itemize}
}

\pf{
    We first prove the identities for elementary processes and then extend them by density.

    For (i), let $H = Z\mathsf{1}_{(U,V]}$. Then
    $$
        (H \cdot M)_t = Z(M_{V \wedge t} - M_{U \wedge t}),
    $$
    so
    $$
        (H \cdot M)^T_t = (H \cdot M)_{t \wedge T}
        = Z(M_{V \wedge T \wedge t} - M_{U \wedge T \wedge t}).
    $$
    On the other hand,
    $$
        H(0,T] = Z\mathsf{1}_{(U \wedge T, V \wedge T]},
    $$
    and therefore
    $$
        (H(0,T]\cdot M)_t
        = Z(M_{V \wedge T \wedge t} - M_{U \wedge T \wedge t}).
    $$
    Similarly,
    $$
        (H \cdot M^T)_t = Z(M^T_{V \wedge t} - M^T_{U \wedge t})
        = Z(M_{V \wedge T \wedge t} - M_{U \wedge T \wedge t}).
    $$
    Hence (i) holds for elementary processes, and then for all $H \in L^2(M)$ by continuity of the integral map.

    For (ii), set
    $$
        N_t = (H \cdot M)_t^2 - \int_0^t H_s^2 \diff [M]_s.
    $$
    Since
    $$
        \mathbb{E}\left[\sup_{t \geq 0}(H \cdot M)_t^2\right]
        \leq 4\mathbb{E}\left[\int_0^\infty H_s^2 \diff [M]_s\right] < \infty,
    $$
    the family $\{N_t : t \geq 0\}$ is uniformly integrable. It therefore suffices to show that $\mathbb{E}N_T = 0$ for every stopping time $T$.

    By part (i),
    $$
        (H \cdot M)_T = (H(0,T]\cdot M)_\infty.
    $$
    Applying the isometry to the integrand $H(0,T]$, we obtain
    $$
    \begin{aligned}
        \mathbb{E}(H \cdot M)_T^2
        &= \mathbb{E}(H(0,T]\cdot M)_\infty^2 \\
        &= \mathbb{E}\int_0^\infty H(0,T]_s^2 \diff [M]_s \\
        &= \mathbb{E}\int_0^T H_s^2 \diff [M]_s.
    \end{aligned}
    $$
    Hence $\mathbb{E}N_T = 0$, and (ii) follows.

    For (iii) First take $H = Z_1 \mathsf{1}_{(S_1,T_1]}$ and $K = Z_2 \mathsf{1}_{(S_2,T_2]}$. Then
    $$
        HK = Z_1Z_2\,\mathsf{1}_{(S,T]},
    $$
    where
    $$
        S = S_1 \vee S_2,
        \qquad
        T = (T_1 \wedge T_2) \vee (S_1 \vee S_2).
    $$
    Hence
    $$
        ((HK)\cdot M)_\infty = Z_1Z_2(M_T - M_S).
    $$
    On the other hand,
    $$
    \begin{aligned}
        (H \cdot (K \cdot M))_\infty
        &= Z_1\Bigl((K \cdot M)_{T_1} - (K \cdot M)_{S_1}\Bigr) \\
        &= Z_1\Bigl((K \cdot M^{T_1})_\infty - (K \cdot M^{S_1})_\infty\Bigr) \\
        &= Z_1Z_2\Bigl[(M_{T_2 \wedge T_1} - M_{S_2 \wedge T_1}) - (M_{T_2 \wedge S_1} - M_{S_2 \wedge S_1})\Bigr] \\
        &= Z_1Z_2(M_T - M_S).
    \end{aligned}
    $$
    Therefore, for every $H,K \in b\varepsilon$,
    $$
        H \cdot (K \cdot M) = (HK)\cdot M.
    $$
    The identity
    $$
        K \cdot (H \cdot M) = (HK)\cdot M
    $$
    is proved in the same way.

    Now let $K \in b\mathcal{P}$. Choose $(K^n)_{n\geq 1} \subset b\varepsilon$ such that
    $$
        K^n \xlongrightarrow{L^2(M)} K.
    $$
    Fix $H \in b\varepsilon$. By the isometry,
    $$
    \begin{aligned}
        \mathbb{E}\left[\left(H \cdot (K\cdot M)-H\cdot(K^n\cdot M)\right)_\infty^2\right]
        &= \mathbb{E}\left[\left(H\cdot\bigl((K-K^n)\cdot M\bigr)\right)_\infty^2\right] \\
        &= \mathbb{E}\left[\int_0^\infty H_s^2 \diff [ (K-K^n)\cdot M ]_s\right] \\
        &= \mathbb{E}\left[\int_0^\infty H_s^2 (K_s-K_s^n)^2 \diff [M]_s\right] \\
        &\leq \norm{H}_\infty^2\,\mathbb{E}\left[\int_0^\infty (K_s-K_s^n)^2 \diff [M]_s\right]
        \to 0.
    \end{aligned}
    $$
    Hence
    $$
        H \cdot (K^n\cdot M) \xlongrightarrow{L^2(\mathcal{F}_\infty)} H \cdot (K\cdot M).
    $$

    On the other hand, for each $n$ we already know that
    $$
        H \cdot (K^n\cdot M) = (HK^n)\cdot M.
    $$
    Since $H \in b\varepsilon$, we also have
    $$
        \mathbb{E}\left[\int_0^\infty (H_sK_s^n-H_sK_s)^2 \diff [M]_s\right]
        \leq \norm{H}_\infty^2\,\mathbb{E}\left[\int_0^\infty (K_s^n-K_s)^2 \diff [M]_s\right]
        \to 0,
    $$
    so $HK^n \xlongrightarrow{L^2(M)} HK$. Applying the isometry once again,
    $$
        (HK^n)\cdot M \xlongrightarrow{L^2(\mathcal{F}_\infty)} (HK)\cdot M.
    $$
    By uniqueness of the $L^2(\mathcal{F}_\infty)$ limit,
    $$
        H \cdot (K\cdot M) = (HK)\cdot M.
    $$
    (iv) First let $H = Z\mathsf{1}_{(S,T]}$. Then
    $$
        (H \cdot M)_t = Z(M_{t \wedge T} - M_{t \wedge S}),
    $$
    so
    $$
        \Delta(H \cdot M)_t = Z\,\mathsf{1}_{(S,T]}(t)\,\Delta M_t = H_t \Delta M_t.
    $$
    Hence the identity holds for all $H \in b\varepsilon$.

    Now let $H \in L^2(M)$, and choose $H^n \in b\varepsilon$ such that
    $$
        H^n \xlongrightarrow{L^2(M)} H.
    $$
    By the isometry,
    $$
        H^n \cdot M \xlongrightarrow{L^2(\mathcal{F}_\infty)} H \cdot M,
    $$
    and hence, after passing to a subsequence,
    $$
        \sup_{t \geq 0}\abs{(H^{n_k}\cdot M)_t - (H\cdot M)_t} \xlongrightarrow{a.s.} 0.
    $$
    Also,
    $$
        \int_0^\infty (H^{n_k}_s - H_s)^2 \diff [M]_s \xlongrightarrow{a.s.} 0
    $$
    along a further subsequence.

    Fix $t \geq 0$. Since
    $$
        \Delta(H^{n_k}\cdot M)_t = H^{n_k}_t \Delta M_t,
    $$
    we distinguish two cases.

    \textbf{Case 1.} If $\Delta M_t = 0$, then $\Delta(H^{n_k}\cdot M)_t = 0$ for every $k$, and therefore
    $$
        \Delta(H \cdot M)_t = 0 = H_t \Delta M_t.
    $$

    \textbf{Case 2.} If $\Delta M_t \neq 0$, then $\Delta[M]_t = (\Delta M_t)^2 > 0$. Hence
    $$
        \int_0^\infty (H^{n_k}_s - H_s)^2 \diff [M]_s
        \geq
        (H^{n_k}_t - H_t)^2 \Delta[M]_t,
    $$
    and so $H^{n_k}_t \to H_t$. Letting $k \to \infty$ in
    $$
        \Delta(H^{n_k}\cdot M)_t = H^{n_k}_t \Delta M_t,
    $$
    we obtain
    $$
        \Delta(H \cdot M)_t = H_t \Delta M_t.
    $$
}
